\clearpage
\section{Reference allocation with fresh store identity}
\label{sec:copying-allocation}
The first concrete store computation is now formalized: reference creation,
Figure~19(a) of the paper (PDF p.~38). The local rule retains its stable-value
premise and the canonical atom tree $D^{\diamond}_{1,1}$. Its actual graph
construction has proved freshness, preservation of the argument, exact store
resolution, renaming compatibility, and forward padding preservation. A
search/allocation subsystem executes the unchanged Figure~16 translation of
$\mathtt{ref}\ n$ in five transitions for every natural literal $n$.

This is an operational extension beyond the preceding search-only milestone.
It is still a subsystem: it omits copying, dereference, assignment, atom equality,
and the other compute rules. Its allocation transitions are rooted local
instances, optionally followed by an unchanged serial continuation. A complete
contextual matching relation, focused-graph correspondence, determinism and
refocusing for the combined machine remain to be proved. No observational
execution instance or Proposition~12.2 theorem is asserted here.

\Needspace{12\baselineskip}
\subsection{The literal allocation pair and its stable premise}

\termdef{copying-allocation-rule}
For a stable graph $G:\star\Rightarrow O$, the redex is serial composition
of the $\mathtt{ref}$ corolla with $G$. The allocated prefix is serial
composition of an instance, $D^{\diamond}_{1,1}$, and a store. It has type
$\star\Rightarrow\star$. Appending $G$ gives the reduct with the same external
interface as the redex. The canonical tree is the original tree constructor;
its weakening and contraction edges are not removed as an optimisation.
All carriers and incidence equations are those of actual finite typed graphs.

The local rule is generated by this pair with a proof of stability of $G$.
Its inversion theorem states exactly that a related pair has such a stable
witness, with equalities to the constructed redex and reduct. The forward
proof eliminates the rule constructor; the converse substitutes the two
equalities and applies that constructor. There is no premise assuming the
resulting programs equivalent or the machine correct.

\lean{Copying/Allocation}{10}{11}{Allocation.before}

\lean{Copying/Allocation}{14}{16}{Allocation.cellPrefix}

\lean{Copying/Allocation}{19}{20}{Allocation.after}

\lean{Copying/Allocation}{23}{24}{Allocation.Rule}

\lean{Copying/Allocation}{27}{31}{Allocation.rule_iff}

\Needspace{12\baselineskip}
\subsection{Fresh allocation and exact store resolution}

\termdef{copying-allocation-fresh-store}
The instance, contraction, weakening and store are explicit edge coordinates
in the nested serial assembly. Every original edge of $G$ is injected into
the other component of the outer serial assembly. Injectivity follows from
the tagged carrier representation. The fresh store is in the prefix component,
so it cannot equal any retained argument edge. Original labels are unchanged,
and the designated instance and store have their required primitive labels.
Freshness is therefore a property of disjoint graph occurrences, independent
of any numerical choice of heap addresses.

\lean{Copying/Allocation}{34}{35}{Allocation.instanceEdge}

\lean{Copying/Allocation}{37}{38}{Allocation.contraction}

\lean{Copying/Allocation}{40}{41}{Allocation.weakening}

\lean{Copying/Allocation}{43}{44}{Allocation.store}

\lean{Copying/Allocation}{47}{48}{Allocation.oldEdge}

\lean{Copying/Allocation}{50}{52}{Allocation.oldEdge_injective}

\lean{Copying/Allocation}{54}{55}{Allocation.store_fresh}

\lean{Copying/Allocation}{57}{58}{Allocation.old_label}

\lean{Copying/Allocation}{60}{61}{Allocation.instance_label}

\lean{Copying/Allocation}{63}{64}{Allocation.store_label}

The instance target is the input of the canonical atom contraction. The
contraction target is the new store input. These two adjacency witnesses,
together with their exact primitive labels, give a two-hop atom trace. Its
endpoint has the store label, so it is a successful resolution. The previously
proved functional atom-tracing theorem implies that any successful trace
starting at this instance has that same store and exactly two hops.
After allocation resets focus to search, the instance rule returns backtrack
at the external input in one step. This last result is about search, whereas
the two-hop trace is the separate store-matching relation; their costs must
not be added together as machine transitions.

\lean{Copying/Allocation}{67}{75}{Allocation.instance_hop}

\lean{Copying/Allocation}{78}{86}{Allocation.contraction_hop}

\lean{Copying/Allocation}{89}{91}{Allocation.resolves}

\lean{Copying/Allocation}{94}{96}{Allocation.resolves_exact}

\lean{Copying/Allocation}{99}{103}{Allocation.returns}

\Needspace{12\baselineskip}
\subsection{Renaming and forward padding preservation}

A recursive isomorphism between two argument graphs extends through the
redex by identity on its $\mathtt{ref}$ head, and through the reduct by identity
on its allocated prefix. Stability transports along the same argument
isomorphism, yielding the corresponding raw allocation rule. Thus the stable
premise is preserved, rather than silently discarded by a representation change.

One recursive padding step leaves the canonical atom tree unchanged except
for its allowance index. The prefix padding comparison follows by applying
serial-padding comparison twice, then the primitive and canonical-tree padding
equalities. The redex comparison uses one serial-padding comparison and the
primitive equality. The reduct comparison uses serial padding and the prefix
comparison. Stability under padding supplies a valid raw allocation at the
higher allowance.

\lean{Copying/AllocationTransport}{9}{10}{Allocation.before_congr}

\lean{Copying/AllocationTransport}{13}{14}{Allocation.after_congr}

\lean{Copying/AllocationTransport}{17}{18}{Allocation.renamed_rule}

\lean{Copying/AllocationTransport}{21}{23}{Allocation.raise_tree}

\lean{Copying/AllocationTransport}{25}{30}{Allocation.prefix_raise}

\lean{Copying/AllocationTransport}{32}{35}{Allocation.before_raise}

\lean{Copying/AllocationTransport}{37}{39}{Allocation.after_raise}

\lean{Copying/AllocationTransport}{42}{44}{Allocation.raised_rule}

\termdef{copying-allocation-up-to-iso}
An allocation up to isomorphism is witnessed by a stable argument and recursive
isomorphisms from its two given graphs to the literal redex and reduct.
Every raw rule supplies identity isomorphisms. Further source and target
renamings compose inversely with the witnesses, proving representation
congruence. To pad such a rule, pad the argument and both witness isomorphisms,
then compose with the redex and reduct padding comparisons. Induction gives
preservation under any finite number of raises.

These are forward operational padding results for allocation pairs. They do
not yet reflect all possible matches from a higher allowance, or prove the
full execution-padding interface. In particular, no padding property is
inferred solely from determinism.

\lean{Copying/AllocationTransport}{47}{48}{Allocation.UpToIso}

\lean{Copying/AllocationTransport}{50}{52}{Allocation.rule_upToIso}

\lean{Copying/AllocationTransport}{54}{57}{Allocation.upToIso_congr}

\lean{Copying/AllocationTransport}{60}{64}{Allocation.upToIso_raise}

\lean{Copying/AllocationTransport}{67}{71}{Allocation.upToIso_pad}

\Needspace{12\baselineskip}
\subsection{Preserved contents and non-aliasing in larger graphs}

An atom hop inside an inserted graph remains a hop after shallow plugging:
map its shared vertex along the inserted-vertex map and use the mapped source
and target memberships. Labels are unchanged. Mapping a finite trace proves
resolution at the same hop count. Comparing that known resolution with any
other successful trace in the filled host, functional terminal-trace uniqueness
forces the same inserted store and count. It follows that the inserted reference
cannot resolve to a retained host edge, because the inserted and retained edge
tags differ. No acyclicity of the host is assumed; the claim applies equally
when the host contains cyclic store structures. It is a resolution-locality
result, not a completed full-machine specimen law.

\lean{Copying/AllocationLocality}{8}{14}{StoreTrace.hop_insert}

\lean{Copying/AllocationLocality}{17}{21}{StoreTrace.resolves_insert}

\lean{Copying/AllocationLocality}{25}{30}{StoreTrace.resolves_insert_exact}

\lean{Copying/AllocationLocality}{33}{40}{StoreTrace.inserted_not_host}

The old-vertex map is the right attachment map of serial composition.
Every old edge retains its source and target lists mapped along this function.
The new store's target list is precisely the singleton containing the mapped
old input. Thus allocation keeps the argument graph and connects the store
to that argument, rather than to an unspecified value. Exact resolution and
freshness exclude successful lookup at any old argument edge. Two independently
allocated graphs placed in parallel also have distinct store occurrences,
because their parallel-component tags differ.

\lean{Copying/AllocationLocality}{47}{48}{Allocation.oldVertex}

\lean{Copying/AllocationLocality}{50}{51}{Allocation.old_source}

\lean{Copying/AllocationLocality}{53}{54}{Allocation.old_target}

\lean{Copying/AllocationLocality}{57}{59}{Allocation.stored_input}

\lean{Copying/AllocationLocality}{62}{65}{Allocation.not_old_store}

\lean{Copying/AllocationLocality}{68}{70}{Allocation.parallel_stores_distinct}

\Needspace{12\baselineskip}
\subsection{Safety and the scope of stability}

For any primitive predicate $P$, the allocation prefix satisfies $P$
whenever $P$ holds of instance, atom weakening, atom contraction and store.
Prove this by serial assembly of the primitive corollas and case analysis on
the canonical tree's edge label. Appending an argument satisfying $P$ preserves
that global property. Specialising to exclusion of focus labels and exclusion
of hole labels proves preservation of focus freedom and hole freedom.
A stable argument supplies both properties, so its allocated graph has both.

The whole allocated graph is not stable in the paper's sense: its designated
store edge would have to have a stable primitive label, but stores are not
stable primitives. Injectivity of primitive labels turns this into a
contradiction. The rule's stability condition concerns its preserved argument,
not the heap containing that argument. This distinction prevents an incorrect
whole-graph stability invariant from being assumed for future store rules.

\lean{Copying/AllocationSafety}{9}{23}{Allocation.prefix_labels}

\lean{Copying/AllocationSafety}{25}{28}{Allocation.after_labels}

\lean{Copying/AllocationSafety}{31}{33}{Allocation.after_focusFree}

\lean{Copying/AllocationSafety}{35}{37}{Allocation.after_holeFree}

\lean{Copying/AllocationSafety}{40}{42}{Allocation.stable_allocation_safe}

\lean{Copying/AllocationSafety}{45}{50}{Allocation.after_not_stable}

\Needspace{12\baselineskip}
\subsection{A graph-changing subsystem and the original translated execution}

\termdef{copying-allocation-subsystem}
A state is a dependent pair of an actual graph and a token coordinate in that
graph. The subsystem transition relation has three constructors: an intrinsic
search step on a fixed graph; a rooted allocation of a stable argument from
rewrite mode to search mode; and that allocation with an unchanged serial
continuation. Both allocation constructors keep the stable premise. The
continuation constructor retains its whole graph, including the original
translation's single-use distributors.

Successful finality is backtrack mode at an external input. Search has no
step there by its proved input obstruction; allocation requires rewrite mode,
so cannot apply. These cases prove final normality and supply a genuine
abstract-machine instance for the subsystem. Determinism of this graph-changing
relation is not established by this construction and is not assumed in the
regression executions below. An allocation followed by any serial continuation
returns its new instance in one search step: use the actual retained instance
edge and its singleton source at the overall input.

\lean{Copying/AllocationExecution}{10}{10}{Allocation.State}

\lean{Copying/AllocationExecution}{14}{20}{Allocation.Step}

\lean{Copying/AllocationExecution}{23}{34}{Allocation.completed_normal}

\lean{Copying/AllocationExecution}{37}{40}{Allocation.machine}

\lean{Copying/AllocationExecution}{43}{48}{Allocation.returns_suffix}

The constant argument is the unchanged translation of the natural literal
$n$ in the empty environment. Its two zero-resource distributor components have
no edges; the remaining constant edge is stable. Eliminating those empty
components proves the global stable-label criterion. The reference graph is
the unchanged translation of $\mathtt{ref}\ n$. The allocated graph replaces
its local reference prefix while retaining the same single-use environment
distributors. Head and value edges are their actual finite assembly coordinates,
and the argument vertex is the head's target port.

\lean{Copying/AllocationExecution}{51}{51}{Allocation.constantArgument}

\lean{Copying/AllocationExecution}{53}{62}{Allocation.constantArgument_stable}

\lean{Copying/AllocationExecution}{65}{65}{Allocation.referenceGraph}

\lean{Copying/AllocationExecution}{68}{68}{Allocation.allocatedGraph}

\lean{Copying/AllocationExecution}{70}{70}{Allocation.referenceHead}

\lean{Copying/AllocationExecution}{71}{71}{Allocation.referenceValue}

\lean{Copying/AllocationExecution}{72}{73}{Allocation.referenceArgument}

The first search interaction enters the reference operation and searches
its argument. The constant returns backtrack at that argument vertex. Returning
through the unary operation then changes to rewrite mode at the root. Each
proof uses the actual edge, label and ordered incidence list. Append the
stable allocation-with-continuation step and the instance-return step to obtain
an exactly five-transition machine run. Its endpoint satisfies successful
finality. This is existence of the displayed five-step execution, not a theorem
that every subsystem execution has that length.

The use of empty environments here is a concrete semantic regression. It does
not restrict the general local allocation theorem, and it does not replace the
required arbitrary-environment Proposition~12.2 endpoint.

\lean{Copying/AllocationExecution}{76}{80}{Allocation.reference_enter}

\lean{Copying/AllocationExecution}{83}{87}{Allocation.reference_value}

\lean{Copying/AllocationExecution}{90}{92}{Allocation.reference_ready}

\lean{Copying/AllocationExecution}{95}{101}{Allocation.reference_run}

\lean{Copying/AllocationExecution}{104}{105}{Allocation.reference_completes}

\Needspace{12\baselineskip}
\subsection{Checking the returned reference and its stored literal}

An atom hop in the left component of serial composition is preserved by
mapping its shared vertex through the left attachment map. Labels are unchanged;
induction transports exact resolution traces. This lifts the local allocation
resolution to the actual post-execution graph with its untouched distributors.
The allocated instance and store names refer to these retained left-component
occurrences. The resulting reference resolves to that store in two atom hops;
any alternative successful resolution agrees on both store and count.

\lean{Copying/AllocationStoreExamples}{9}{17}{StoreTrace.hop_serial_left}

\lean{Copying/AllocationStoreExamples}{19}{22}{StoreTrace.resolves_serial_left}

\lean{Copying/AllocationStoreExamples}{28}{29}{Allocation.allocatedInstance}

\lean{Copying/AllocationStoreExamples}{31}{32}{Allocation.allocatedStore}

\lean{Copying/AllocationStoreExamples}{35}{37}{Allocation.allocated_resolves}

\lean{Copying/AllocationStoreExamples}{39}{41}{Allocation.allocated_resolves_exact}

The allocated value occurrence is the original constant edge retained
through both serial assemblies. Its label is the literal $n$, and the target
list of the resolved store equals the source list of this value edge, by the
actual incidence definitions. Instantiating the execution and resolution
results at $n=1$ gives the requested example's allocation subexpression:
a five-transition run, a fresh reference resolving to its store, and that
store attached to the original literal one. Final normality rules out further
subsystem transitions after the reference has returned.

This does not yet run the enclosing let, copy or invoke the lambda, or execute
a dereference. Those operations and their template proof obligations remain
on the route to Proposition~12.2.

\lean{Copying/AllocationStoreExamples}{44}{45}{Allocation.allocatedValue}

\lean{Copying/AllocationStoreExamples}{47}{48}{Allocation.allocated_value_label}

\lean{Copying/AllocationStoreExamples}{51}{52}{Allocation.allocated_store_value}

\lean{Copying/AllocationStoreExamples}{55}{60}{Allocation.ref_one}

\lean{Copying/AllocationStoreExamples}{63}{64}{Allocation.reference_final_normal}

Validation checks all 304 build jobs and all three axiom audits: six
abstract, 477 graph and 460 copying declarations. Only standard Lean logical
axioms occur; there are no project-specific axioms or unfinished proofs.
The estimated completion of Proposition~12.2 is roughly 55--60\% by mathematical
effort, with a broad 40--65\% uncertainty range (26 October 2026). This is a
judgement about the remaining mathematics, not a declaration count. The
largest unknowns remain copying/matching, combined-machine laws, and the
Table~3 template robustness proofs.

